% Part: proof-theory
% Chapter: proof-search
% Section: search-algorithm

\documentclass[../../../include/open-logic-section]{subfiles}

\begin{document}

\olfileid{pt}{ps}{sal}

\olsection{Search Algorithm for \Log{G3c}}

We describe a search algorithm for \Log{G3c}. It is not the only
possible one, and not even the most efficient one. But it is correct:
if a sequent $\Gamma \Sequent \Delta$ has !!a{proof}, it will
eventually halt and find !!a{proof}. It is also fool-proof: if it
halts without having found !!a{proof}, or never halts, then $\Gamma
\Sequent \Delta$ is not valid to begin with.

A key necessary feature of the algorithm is that it ensures that every
!!{formula} in $\Gamma \Sequent \Delta$, or in a sequent produced in
the proof search, is ``reduced,'' i.e., it is treated as the principal
formula of an inference rule applied backwards, and that this happens
as often as necessary. We call this feature \emph{fairness}

The algorithm works in stages. The output of each stage is a tree of
sequents, and the output of the next stage is produced by reducing
!!a{formula} in every topmost sequent that isn't already an axiom.

At stage $0$ we start with $\Gamma \Sequent \Delta$. In order to
ensure fairness we assign numbers as indices to the !!{formula}s in
$\Gamma$ and $\Delta$ in such a way that different !!{formula}s
receive different numbers. E.g., think of counting off the
!!{formula}s from left to right. We write the numbers as superscripts.
E.g., if we are searching for !!a{proof} of $!A, !B \Sequent !C$, the
output of stage~$0$ is $!A^1, !B^2 \Sequent !C^3$. The largest~$n$
occurring as a superscript in such an indexed sequent is its
\emph{maximal index} ($3$~in the example).

If the sequent contains quantifiers, we will also need to know which
terms to use when reducing $\lexists[x][!A(x)]$ on the right or
$\lforall[x][!A(x)]$ on the left. We will only need terms built from
the !!{constant}s $\Obj c_0$, $\Obj c_1$, $\Obj c_2$, \dots, using any
of the function symbols occurring in~$\Gamma$ and $\Delta$. We assume
that the set of all these terms is given in some fixed enumeration, so
that we can talk about, e.g., ``the first !!{constant} not occuring in
$\Gamma \Sequent \Delta$.'' Each indexed sequent in the tree the
algorithm produces has an associated set of !!{constant}s. The set of
!!{constant}s assigned to the sequent in stage~$0$ is the set of all
!!{constant}s occuring in it (if there are any), and $\{\Obj c_0\}$ if
it contains no !!{constant}s.

Suppose the algorithm has generated a tree of indexed sequents, with
associated sets of !!{constant}s, at stage~$n$. At stage~$n+1$, we do
the following, for each topmost sequent. 

If there is !!a{formula}~$!A$ that occurs on both the left and the
right of the sequent, or the left contains~$\lfalse$, we do nothing:
such sequents have !!{proof}s in~\Log{G3c}, and for this particular
topmost sequent there is nothing left to do.

If there is no non-atomic !!{formula} in the sequent, then abort the
algorithm. There is no proof of $\Gamma \Sequent \Delta$.

Otherwise, pick the non-atomic !!{formula}~$!A$ with the smallest
index~$i$. Then the topmost sequent in question is either $!A^i, \Pi
\Sequent \Lambda$ or $\Pi \Sequent \Lambda, !A^i$. Let $k$ be the
maximum index of the sequent, and $C$ the set of !!{constant}s
assigned to it.

We ``reduce'' $!A^i$, i.e., we generate new topmost sequents according
to the inference rule determined by the !!{main operator} of~$!A$
and according to whether $!A^i$ appears on the left or the right.

\begin{enumerate}
  \item $!A \ident !B \land !C$ and occurs on the left: Add one new
  topmost sequent above $(!B \land !C)^i, \Pi \Sequent \Lambda$:
  \[
  \Axiom$!B^k, !C^{k+1}, \Pi \fCenter \Lambda$
  \RightLabel{\LeftR{\land}}
  \UnaryInf$(!B \land !C)^i, \Pi \fCenter \Lambda$
  \DisplayProof
  \]
  The set of !!{constant}s assigned to the new sequent is~$C$.
  \item $!A \ident !B \land !C$ and occurs on the right: Add two new
  topmost sequents above $\Pi \Sequent \Lambda, (!B \land !C)^i$:
  \[
  \Axiom$\Pi \fCenter \Lambda, !B^k$
  \Axiom$\Pi \fCenter \Lambda, !C^k$
  \RightLabel{\RightR{\land}}
  \BinaryInf$\Pi \fCenter \Lambda, (!B \land !C)^i$
  \DisplayProof
  \]
  The set of !!{constant}s assigned to the new sequent is~$C$.

  \item The cases where $!A \ident \lnot !B$, $!A \ident !B \lor !C$,
  and $!A \ident !B \lif !C$ are similar and left as exercises.

  \item $!A \ident \lexists[x][!B(x)]$ and occurs on the left: Add one new
  topmost sequent above $\lexists[x][!B(x)]^i, \Pi \Sequent \Lambda$:
  \[
  \Axiom$!B(c)^k, \Pi \fCenter \Lambda$
  \RightLabel{\LeftR{\lexists}}
  \UnaryInf$\lexists[x][!B(x)]^i, \Pi \fCenter \Lambda$
  \DisplayProof
  \]
  where $c$ is the first !!{constant} not in~$C$. The set of
  !!{constant}s assigned to the new sequent is $C \cup \{c\}$.
  \item $!A \ident \lexists[x][!B(x)]$ and occurs on the right: Add one new
  topmost sequent above $\Pi \Sequent \Lambda, \lexists[x][!B(x)]$:
  \[
  \Axiom$\Pi \fCenter \Lambda, \lexists[x][!B(x)]^k, !B(t)^{k+1}$
  \RightLabel{\LeftR{\lexists}}
  \UnaryInf$\Pi \fCenter \Lambda, \lexists[x][!B(x)]^i$
  \DisplayProof
  \]
  where $t$ is the first term containing only !!{constant}s in~$C$ not
  already used in the reduction of $\lexists[x][!B(x)]$ on the right
  below this sequent (or $\Obj c_0$ if all such terms have been used).
  The set of !!{constant}s assigned to the new sequent is $C$.
  \item The cases where $!A \ident \lforall[x][!B(x)]$ are similar and
  left as exercises.
\end{enumerate}

If at stage $n+1$ we have not added a new sequent to any topmost
sequent, the algorithm terminates successfully. In this case we have
found !!a{proof} of~$\Gamma \Sequent \Delta$, obtained from the tree
at stage $n+1$ by erasing all indices. The topmost sequents are all of
the form $!A, \Pi \Sequent \Lambda, !A$ or of the form $\lfalse, \Pi
\Sequent \Lambda$. All inferences are correct inferences of~\Log{G3c}.
This is obvious for the propositional inferences. Observe that at each
step, the set~$C$ of !!{constant}s assigned to a sequent contains all
the !!{constant}s occuring in the sequent. Hence, in the reduction
using \LeftR{\lexists} and \RightR{\lforall}, the eigenvariable
condition is satisfied, since the eigenvariable~$c$ was !!a{constant}
not in the set~$C$ assigned to the conclusion, $c$ also did not occur
in the conclusion. We have:

\begin{prop}
The proof search algorithm for \Log{G3c} is correct: if it terminates
successfully, the starting sequent~$\Gamma \Sequent \Delta$ has
!!a{proof}.
\end{prop}





\end{document}
